\documentclass[11pt,a4paper]{article}
\usepackage[T1]{fontenc}
\usepackage{lmodern}
\usepackage[margin=27mm]{geometry}
\usepackage{amsmath,amssymb,amsthm,mathtools}
\usepackage{microtype}
\usepackage{booktabs,array}
\usepackage{enumitem}
\usepackage[hidelinks]{hyperref}
\setlength{\parskip}{5pt plus 1pt}
\setlength{\parindent}{0pt}
\setlength{\emergencystretch}{3em}
\setlist{itemsep=2pt,topsep=3pt}
\raggedbottom
\newtheorem{theorem}{Theorem}[section]
\newtheorem{lemma}[theorem]{Lemma}
\newtheorem{corollary}[theorem]{Corollary}
\newtheorem{proposition}[theorem]{Proposition}
\theoremstyle{definition}
\newtheorem{definition}[theorem]{Definition}
\newtheorem{example}[theorem]{Example}
\theoremstyle{remark}
\newtheorem{remark}[theorem]{Remark}
\newcommand{\Z}{\mathbb{Z}}
\newcommand{\N}{\mathbb{N}}
\newcommand{\Q}{\mathbb{Q}}
\newcommand{\F}{\mathbb{F}}
\newcommand{\ps}[1]{#1[[x]]}
\newcommand{\coeff}[2]{[x^{#1}]\,#2}
\newcommand{\hl}[1]{\textbf{#1}}
\newcommand{\para}[1]{\par\medskip\noindent\textbf{#1}\enspace}
\hypersetup{pdftitle={The remainders of the Catalan-Stirling numbers A052895 are eventually periodic},pdfauthor={Dev D. Goyal},pdfsubject={OEIS A052895}}
\title{The remainders of the Catalan--Stirling numbers A052895 are eventually periodic\\[6pt]\large A short proof of a known periodicity theorem, with a formal verification}
\author{Dev D. Goyal}
\date{October 2026}
\begin{document}
\maketitle
\begin{abstract}
The sequence A052895 of the On-Line Encyclopedia of Integer Sequences (OEIS) has exponential generating function $C(e^x-1)$, where $C(y)=\sum_kC_ky^k$ is the generating function of the Catalan numbers $C_k=1,1,2,5,14,\dots$; equivalently $a(n)=\sum_kC_k\,k!\,S(n,k)$. Peter Bala conjectured that, for every whole number $m\ge1$, the remainders of $a(n)$ on
division by $m$ eventually repeat in a cycle whose length divides $\varphi(m)$, Euler's totient of $m$. This is a special
case of a 2026 theorem of Payam Danesh. We give a short, self-contained proof that works for every sequence of the form
$\sum_k c_k\,k!\,S(n,k)$ with integer weights $c_k$, where $S(n,k)$ are the Stirling numbers of the second kind: for all
$n\ge m$, $a(n+\varphi(m))\equiv a(n)\pmod m$. The proof uses only Euler's theorem and the fact that $m$ divides $k!$
when $k\ge m$. All results are also proved formally in the Lean proof assistant.
\end{abstract}
\begin{center}\begin{minipage}{0.88\linewidth}\small
\textbf{How this paper was produced.} This paper reports the result of an autonomous AI research effort.
ChatGPT (OpenAI), using the GPT-6 model with Pro reasoning and working from prompts by Dev~D.~Goyal, wrote the self-contained proof of this known result, its Lean~4 formalization, and a first draft of the paper.
Dev~D.~Goyal chose the problems and asked for checks and revisions, but did not supply proof ideas.
Claude Opus~5.5, used with Claude Code, then checked the work independently: it rebuilt the formal proofs from
source and confirmed that they rely on no unproved assumptions, recomputed the sequence from its definition and tested
every claim against the published data, and read every step of the proofs. It also rewrote this paper for clarity.
Dev~D.~Goyal has verified the results and is responsible for them. The paper has not been peer reviewed.
The section ``About this paper'' at the end gives the details.
\end{minipage}\end{center}

\section{Introduction}

\subsection{The sequence}\label{sec:sequence}
The \emph{On-Line Encyclopedia of Integer Sequences} (OEIS)~\cite{oeisdb} is a large public database of integer sequences.
Besides the terms, many entries record formulas, interpretations, and \emph{conjectures}: patterns observed in the data
that the contributor could not prove. This paper is about entry A052895~\cite{oeis}, which reads:
\begin{quote}\small\ttfamily\raggedright Expansion of e.g.f. (1/2)/(exp(x) - 1) * (1 - (5 - 4*exp(x))\^{}(1/2)).\par\smallskip\hfill\textrm{--- OEIS A052895, INRIA Encyclopedia of Combinatorial Structures (2000)}\end{quote}

Here \texttt{exp(x)} is the exponential series $e^x=\sum_{n\ge0}x^n/n!$, and the ``e.g.f.'' (exponential
generating function) of a sequence $a(n)$ is the series $\sum_na(n)\,x^n/n!$.

Like many sequences in combinatorics, this one is a \emph{weighted Stirling transform}. For each $k$, let
$k!\,S(n,k)$ be the number of ways to put $n$ labelled objects into $k$ labelled boxes so that no box is empty
($S(n,k)$, the Stirling number of the second kind, counts the ways when the boxes are unlabelled). Then
\begin{equation}\label{eq:def}
a(n)=\sum_{k=0}^{n}C_k\;k!\;S(n,k),
\end{equation}
where $C_k=\frac1{k+1}\binom{2k}{k}$ is the $k$-th Catalan number ($1,1,2,5,14,42,\dots$). In words: $a(n)$ counts the ways to arrange $n$ labelled objects into a row of
nonempty groups, where an arrangement with $k$ groups is counted $C_k$ times. The two descriptions agree: writing $y=e^x-1$, we have $5-4e^x=1-4y$, so the e.g.f.\ in the entry is
$\frac{1-\sqrt{1-4y}}{2y}=\sum_kC_ky^k$, the generating function of the Catalan numbers, and Lemma~\ref{lem:egf} below
turns this into \eqref{eq:def}. The entry also records \eqref{eq:def} directly (V.~Kruchinin).

The first terms are 1, 1, 5, 43, 545, 9211, 195305, 4990483, 149371745, 5128125451, \dots{} (from $n=0$).

\subsection{The question}
The entry contains the following conjecture.
\begin{quote}\small\ttfamily\raggedright Conjecture: for fixed k = 1,2,..., the sequence a(n) (mod k) is eventually periodic with the exact period dividing phi(k), where phi(k) is the Euler totient function A000010. For example, modulo 10 the sequence becomes (1, 1, 5, 3, 5, 1, 5, 3, 5, ...), with an apparent period 1, 5, 3, 5 of length 4 = phi(10) beginning at a(1).\par\smallskip\hfill\textrm{--- OEIS A052895, Peter Bala, January 2018}\end{quote}

Here $\varphi(m)$, \emph{Euler's totient}, is the number of integers between $1$ and $m$ that have no common factor with
$m$; for example $\varphi(10)=4$ (the numbers $1,3,7,9$) and $\varphi(7)=6$. A sequence of remainders is
\emph{eventually periodic} with period $d$ if, from some index on, it repeats every $d$ steps.
The table shows the remainders modulo $10$, where $\varphi(10)=4$.
The indices in bold are those with $n\ge10$, from which point our theorem guarantees that the pattern repeats
every 4 steps.
\begin{center}\small\setlength{\tabcolsep}{3.2pt}
\begin{tabular}{>{\raggedright\arraybackslash}p{6.4em}>{\centering\arraybackslash}p{1.15em}>{\centering\arraybackslash}p{1.15em}>{\centering\arraybackslash}p{1.15em}>{\centering\arraybackslash}p{1.15em}>{\centering\arraybackslash}p{1.15em}>{\centering\arraybackslash}p{1.15em}>{\centering\arraybackslash}p{1.15em}>{\centering\arraybackslash}p{1.15em}>{\centering\arraybackslash}p{1.15em}>{\centering\arraybackslash}p{1.15em}>{\centering\arraybackslash}p{1.15em}>{\centering\arraybackslash}p{1.15em}>{\centering\arraybackslash}p{1.15em}>{\centering\arraybackslash}p{1.15em}>{\centering\arraybackslash}p{1.15em}>{\centering\arraybackslash}p{1.15em}}
\toprule
$n$ & 0 & 1 & 2 & 3 & 4 & 5 & 6 & 7 & 8 & 9 & \hl{10} & \hl{11} & \hl{12} & \hl{13} & \hl{14} & \hl{15} \\ \midrule
$a(n)\bmod 10$ & 1 & 1 & 5 & 3 & 5 & 1 & 5 & 3 & 5 & 1 & \hl{5} & \hl{3} & \hl{5} & \hl{1} & \hl{5} & \hl{3} \\
\bottomrule
\end{tabular}
\\[4pt]
\begin{tabular}{>{\raggedright\arraybackslash}p{6.4em}>{\centering\arraybackslash}p{1.15em}>{\centering\arraybackslash}p{1.15em}>{\centering\arraybackslash}p{1.15em}>{\centering\arraybackslash}p{1.15em}>{\centering\arraybackslash}p{1.15em}>{\centering\arraybackslash}p{1.15em}>{\centering\arraybackslash}p{1.15em}>{\centering\arraybackslash}p{1.15em}>{\centering\arraybackslash}p{1.15em}>{\centering\arraybackslash}p{1.15em}>{\centering\arraybackslash}p{1.15em}>{\centering\arraybackslash}p{1.15em}>{\centering\arraybackslash}p{1.15em}>{\centering\arraybackslash}p{1.15em}>{\centering\arraybackslash}p{1.15em}>{\centering\arraybackslash}p{1.15em}}
\toprule
$n$ & \hl{16} & \hl{17} & \hl{18} & \hl{19} & \hl{20} & \hl{21} & \hl{22} & \hl{23} & \hl{24} & \hl{25} & \hl{26} & \hl{27} & \hl{28} & \hl{29} & \hl{30} & \hl{31} \\ \midrule
$a(n)\bmod 10$ & \hl{5} & \hl{1} & \hl{5} & \hl{3} & \hl{5} & \hl{1} & \hl{5} & \hl{3} & \hl{5} & \hl{1} & \hl{5} & \hl{3} & \hl{5} & \hl{1} & \hl{5} & \hl{3} \\
\bottomrule
\end{tabular}
\end{center}


\subsection{What this paper proves}
\textbf{Main result} (Theorem~\ref{thm:main} and Corollary~\ref{cor:seq}). \emph{For every integer $m\ge1$ and every
$n\ge m$, $a(n+\varphi(m))\equiv a(n)\pmod m$.} So from $n=m$ on, the remainders modulo $m$ repeat with period
$\varphi(m)$, and the shortest period divides $\varphi(m)$. This proves Bala's conjecture for this sequence.

The proof never uses the particular weights $C_k$: it works for every sequence of the form \eqref{eq:def} with
integer weights. That covers the more general conjecture Bala stated in several entries, for every sequence whose
exponential generating function is $G(e^x-1)$ with $G$ a power series with integer coefficients.

\subsection{What was known}
The result is \textbf{not new}. Payam Danesh~\cite{danesh} proved the general statement in 2026 (Theorem~4.1 and
Corollary~4.2 there), with a stronger conclusion: the period divides Carmichael's function $\lambda(m)$, the smallest
$e\ge1$ such that $a^e\equiv1\pmod m$ for every $a$ with no common factor with $m$. Since $\lambda(m)$ divides
$\varphi(m)$, our statement follows from Danesh's. Earlier, B.~Schreyer~\cite{schreyer} studied the periodicity of the
Fubini numbers, the case where all weights equal $1$. Our contribution is a very short, self-contained proof that a
non-specialist can follow, and a formal verification of it in Lean.

\subsection{Why it is interesting}
Weighted Stirling transforms appear throughout combinatorics, whenever objects are built by first splitting a set into
blocks and then decorating the arrangement of blocks; the OEIS contains many of them. Knowing that their remainders are
periodic is useful in practice: it makes $a(n)\bmod m$ computable for huge $n$, and it gives a strong consistency check
on computed data. Bala noticed the same phenomenon in many entries; a single theorem explains all of them.

\subsection{How to read this paper}
Section~\ref{sec:background} explains the notions used. Section~\ref{sec:results} states the results and
Section~\ref{sec:proofs} proves them. Section~\ref{sec:checks} describes how the results were checked by computer and by a
formal proof, and Section~\ref{sec:discussion} discusses the scope.
\section{Background}\label{sec:background}
This section collects the definitions and facts used later. Readers who know them can skip it.
\subsection{Remainders and congruences}
For an integer $m\ge2$ we write $a\equiv b\pmod m$ (``$a$ is congruent to $b$ modulo $m$'') when $a-b$ is divisible by
$m$, that is, when $a$ and $b$ leave the same remainder after division by $m$. For example, $17\equiv7\pmod{10}$, since both leave remainder $7$. The
remainder of $a$ on division by $m$ is called its \emph{residue} modulo $m$; it is one of $0,1,\dots,m-1$. Saying that
$a$ is odd is the same as saying $a\equiv1\pmod2$.

Residues can be added and multiplied consistently, as on a clock. The residues $0,1,\dots,m-1$ with this ``clock
arithmetic'' form a number system written $\Z/m\Z$. When $m=p$ is a prime we also write $\F_p$. In $\F_2=\{0,1\}$ we
have $1+1=0$, so adding and subtracting are the same thing.

For power series with integer coefficients, $U\equiv V\pmod m$ means that \emph{every} coefficient of $U-V$ is
divisible by $m$. Equivalently, the two series become equal when every coefficient is replaced by its residue; we call
this \emph{reducing} the series modulo $m$. The key fact is simple:

\emph{Reduction modulo $m$ is compatible with all the operations above.} To reduce a sum, a product, a substitution, or
an inverse, we may reduce the ingredients first and then compute with clock arithmetic. This holds because every
coefficient of the result is obtained from finitely many coefficients of the ingredients by adding, subtracting and
multiplying, and these operations are compatible with remainders.

\subsection{Stirling numbers}
The \emph{Stirling number of the second kind} $S(n,k)$ is the number of ways to split a set of $n$ objects into $k$
nonempty groups, when the groups are not labelled. For example $S(3,2)=3$: the splits of $\{1,2,3\}$ into two groups
are $\{1\}\{2,3\}$, $\{2\}\{1,3\}$ and $\{3\}\{1,2\}$. Labelling the $k$ groups multiplies the count by $k!$. Equivalently,
$S(0,0)=1$, $S(n,0)=S(0,k)=0$ for $n,k>0$, and
\[
S(n+1,k+1)=(k+1)\,S(n,k+1)+S(n,k),
\]
because the new object either joins one of the $k+1$ existing groups or forms a group of its own. This recurrence is the
definition used in Lean. Note that $S(n,k)=0$ when $k>n$. A standard reference is~\cite{concrete}.

\begin{lemma}[Powers via Stirling numbers]\label{lem:power}
For all integers $a,n\ge0$,
\[
a^n=\sum_{k=0}^{a}\binom{a}{k}\,k!\,S(n,k).
\]
\end{lemma}

\begin{proof}
Count the functions from an $n$-element set $N$ to an $a$-element set $T$. There are $a^n$ of them. Now sort them by the
size $k$ of their image. A function with an image of size $k$ is described by splitting $N$ into the $k$ nonempty groups
of elements with the same value ($S(n,k)$ ways), and sending the groups to distinct elements of $T$
($a(a-1)\cdots(a-k+1)=\binom ak k!$ ways).
\end{proof}

\subsection{Euler's theorem}
\begin{lemma}\label{lem:euler}
For all integers $a\ge0$, $m\ge1$ and $n\ge m$, we have $a^{\,n+\varphi(m)}\equiv a^n\pmod m$.
\end{lemma}

\begin{proof}
Write $m$ as a product of prime powers $p^e$; it is enough to check the congruence modulo each $p^e$. If $p$ divides
$a$, both sides are divisible by $p^e$, because $n\ge m\ge p^e\ge e$. If $p$ does not divide $a$, Euler's theorem gives
$a^{\varphi(p^e)}\equiv1\pmod{p^e}$. Since $\varphi(p^e)$ divides $\varphi(m)$ ($\varphi$ is multiplicative), also
$a^{\varphi(m)}\equiv1$, and multiplying by $a^n$ gives the claim.
\end{proof}

The condition $n\ge m$ is needed only for the first case. For example, $2^{0+\varphi(4)}=4\not\equiv1=2^0\pmod 4$, but
$2^{n+2}\equiv2^n\equiv0\pmod4$ for all $n\ge2$.

\subsection{Exponential generating functions}
A \emph{formal power series} $\sum_{n\ge0}u_nx^n$ is just its list of coefficients; we never substitute a number for $x$,
so convergence plays no role. Series are added and multiplied like polynomials. The \emph{exponential generating
function} (e.g.f.) of a sequence $f(0),f(1),\dots$ is the series $\sum_nf(n)\,x^n/n!$. The OEIS often defines a sequence by
its e.g.f.; the next lemma connects such definitions with \eqref{eq:def}.

\begin{lemma}\label{lem:egf}
In formal power series with rational coefficients, for every $k\ge0$,
\[
\frac{(e^x-1)^k}{k!}=\sum_{n\ge0}S(n,k)\,\frac{x^n}{n!},
\qquad\text{and hence}\qquad
\sum_{k\ge0}c_k\,(e^x-1)^k=\sum_{n\ge0}\Bigl(\sum_{k=0}^nc_k\,k!\,S(n,k)\Bigr)\frac{x^n}{n!}.
\]
\end{lemma}

\begin{proof}
If $F=\sum f(n)x^n/n!$ and $H=\sum h(n)x^n/n!$, then $FH=\sum_n\bigl(\sum_i\binom ni f(i)h(n-i)\bigr)x^n/n!$. Applying this
$k-1$ times to $e^x-1=\sum_{n\ge1}x^n/n!$, the number $n!\,\coeff{n}{(e^x-1)^k}$ is the sum of the multinomial coefficients
$\frac{n!}{n_1!\cdots n_k!}$ over all $n_1+\dots+n_k=n$ with every $n_i\ge1$. That sum counts the ways to put $n$ labelled
objects into $k$ labelled nonempty boxes, which is $k!\,S(n,k)$. For the second formula, note that $(e^x-1)^k$ has no
terms below $x^k$, so only finitely many $k$ contribute to each coefficient.
\end{proof}


\section{Results}\label{sec:results}

For an arbitrary sequence of integers $c_0,c_1,c_2,\dots$ (no other assumption), put
\begin{equation}\label{eq:transform}
T_c(n)=\sum_{k=0}^{n}c_k\,k!\,S(n,k).
\end{equation}

\begin{theorem}\label{thm:main}
For every integer $m\ge1$ and every $n\ge m$,
\[
T_c\bigl(n+\varphi(m)\bigr)\equiv T_c(n)\pmod m .
\]
Consequently the sequence of remainders $T_c(m),T_c(m+1),T_c(m+2),\dots$ modulo $m$ is periodic, and its shortest period
divides $\varphi(m)$.
\end{theorem}

\begin{corollary}\label{cor:seq}
For the sequence $a(n)$ of OEIS A052895 (the weights $c_k=C_k$) and every $m\ge1$:
$a(n+\varphi(m))\equiv a(n)\pmod m$ for all $n\ge m$. Hence $a(n)\bmod m$ is eventually periodic, with shortest period
dividing $\varphi(m)$. This is Bala's conjecture.
\end{corollary}

\section{Proofs}\label{sec:proofs}

\para{The idea in brief.} Powers $k^n$ repeat modulo $m$ with period $\varphi(m)$ once $n\ge m$ (Euler).
Lemma~\ref{lem:power} writes the powers $k^n$ in terms of the numbers $j!\,S(n,j)$ with $j\le k$, in a triangular way, so
the periodicity can be passed on to each $j!\,S(n,j)$, one $j$ at a time, without dividing. Finally, every term with
$j\ge m$ is divisible by $m$, because $m$ divides $j!$; so only finitely many terms matter modulo $m$.

\begin{proof}[Proof of Theorem~\ref{thm:main}]
If $m=1$ every congruence modulo $1$ holds, so let $m\ge2$, fix $n\ge m$, and write $\Phi=\varphi(m)$.

\emph{Step 1: the terms with $k<m$.} We show, by induction on $k$, that
\[
k!\,S(n+\Phi,k)\equiv k!\,S(n,k)\pmod m\qquad\text{for }0\le k<m.
\]
Apply Lemma~\ref{lem:power} with $a=k$ to the two exponents $n+\Phi$ and $n$:
\[
k^{n+\Phi}=\sum_{j=0}^{k}\binom kj\,j!\,S(n+\Phi,j),\qquad
k^{n}=\sum_{j=0}^{k}\binom kj\,j!\,S(n,j).
\]
The left sides are congruent modulo $m$ by Lemma~\ref{lem:euler}. By the induction hypothesis, the terms with $j<k$ on the
right are congruent too. So the remaining terms, those with $j=k$ (where $\binom kk=1$), are congruent. Note that we never
divide by $k!$, which would not be allowed modulo $m$.

\emph{Step 2: the terms with $k\ge m$.} If $k\ge m$, then $m$ divides $k!$, so these terms are $\equiv0\pmod m$ in both
$T_c(n+\Phi)$ and $T_c(n)$.

\emph{Step 3.} By Steps 1 and 2,
\[
T_c(n+\Phi)\equiv\sum_{k<m}c_k\,k!\,S(n+\Phi,k)\equiv\sum_{k<m}c_k\,k!\,S(n,k)\equiv T_c(n)\pmod m .
\]

\emph{The shortest period.} Let $f(t)$ be the remainder of $T_c(m+t)$ modulo $m$. We have shown that $\Phi$ is a period of
$f$, so $f$ has a shortest period $d\ge1$. If $p$ is any period, divide: $p=qd+r$ with $0\le r<d$. Then
$f(t+r)=f(t+r+qd)=f(t+p)=f(t)$ for every $t$, so $r$ is also a period, and $r=0$ by the minimality of $d$. Hence $d$
divides every period, in particular $\Phi$.
\end{proof}

Corollary~\ref{cor:seq} is the case $c_k=C_k$ of Theorem~\ref{thm:main}, because $a(n)=T_c(n)$ by \eqref{eq:def}.

\section{Checking the results}\label{sec:checks}

A mathematical proof can contain mistakes, so the results were checked in two independent ways.

\para{Numerical check.}
We computed $a(n)$ from the e.g.f.\ given in the OEIS entry, which reproduces all the terms published there, and
checked that $a(n+\varphi(m))\equiv a(n)\pmod m$ for every $m\le60$ and every $n$ with $m\le n<160$.
The computation was done with the PARI/GP computer algebra system, using code that was written separately from the
formal proofs (file \texttt{verification/independent\_checks.gp} in the accompanying repository~\cite{repo}).
A computation can only test finitely many cases, so it is a check on the definitions and on the indexing, not a proof.

\para{Formal proof in Lean.}
Lean~\cite{lean} is a \emph{proof assistant}: a program in which mathematical statements and their proofs are written
in a precise formal language, and every step is checked mechanically down to the basic rules of logic. All results of
this paper are proved in Lean~4 (version 4.19.0) on top of Mathlib~\cite{mathlib}, Lean's library of formalized
mathematics. The statements are collected in the file \texttt{lean/Entries/A052895.lean} of the repository,
in the namespace \texttt{FinalCollection.A052895}. There $a(n)$ is defined by the finite sum \eqref{eq:def}. The general theorem for arbitrary integer weights is \texttt{More.Stirling.transform\_totient\_period}. The link between the sum and the e.g.f.\ in the OEIS (Lemma~\ref{lem:egf}) is proved in this paper but not in Lean.
The table lists where each result of this paper is proved.
\begin{center}\small\renewcommand{\arraystretch}{1.25}
\begin{tabular}{>{\raggedright\arraybackslash}p{0.52\linewidth}>{\raggedright\arraybackslash}p{0.4\linewidth}}
\toprule
Statement in this paper & Lean theorem(s)\\
\midrule
$a(n)$ is the sum \eqref{eq:def} (the definition used in Lean) & \texttt{source\_\allowbreak{}formula} \\
Corollary~\ref{cor:seq}: $a(n+\varphi(m))\equiv a(n)\pmod m$ for $n\ge m$ & \texttt{period} \\
The shortest period of $t\mapsto a(m+t)\bmod m$ divides $\varphi(m)$ & \texttt{least\_\allowbreak{}period\_\allowbreak{}divides} \\
That shortest period exists, is positive, is a period, and is the shortest & \texttt{least\_\allowbreak{}period\_\allowbreak{}specification} \\
\bottomrule
\end{tabular}
\end{center}
The Lean project was rebuilt from source for this paper, and Lean was asked to list the basic assumptions (axioms)
that each of these theorems ultimately depends on. Every one depends only on the three standard axioms of Lean's logic
(\texttt{propext}, \texttt{Classical.choice}, \texttt{Quot.sound}); no step is left unproved. A formal proof certifies
exactly the statement that Lean checks. Whether that statement says what the OEIS entry means is a separate question,
so the Lean definitions were also compared, side by side, with the definition in the OEIS entry
(Section~\ref{sec:sequence}).

\section{Discussion}\label{sec:discussion}

\para{Stronger known results.} Danesh~\cite{danesh} proves that the period divides $\lambda(m)$, which can be much smaller
than $\varphi(m)$ (for example $\lambda(16)=4$ while $\varphi(16)=8$). The Lean formalization proves only the
$\varphi(m)$ version.

\para{Where the periodicity starts.} Our theorem guarantees periodicity from $n=m$ on. This is enough for the conjecture,
but it is not always the true starting point: for instance, modulo $10$ the remainders of this sequence are already periodic from $n=1$. Finding the exact starting point and the
exact period in general is a natural question that this paper does not address.

\para{Which sequences are covered.} The theorem needs integer weights $c_k$. For example, the Bell numbers
$\sum_kS(n,k)$ correspond to $c_k=1/k!$ and are not covered. Their remainders behave differently: modulo $2$ they are
$1,1,0,1,1,0,\dots$, with period $3$, which does not divide $\varphi(2)=1$.


\section*{About this paper}\label{sec:about}
\para{How the work was done.}
This paper is one of 34 papers on conjectures from the OEIS that came out of an autonomous AI research effort directed
by Dev~D.~Goyal~\cite{repo}. The roles were as follows.
\begin{itemize}
\item \emph{Dev~D.~Goyal} started the project, chose which problems to pursue, asked for proofs, formal verification and
revisions, and decided what to publish. Dev~D.~Goyal did not supply proof strategies or intermediate arguments.
\item \emph{ChatGPT} (OpenAI), using the GPT-6 model with Pro reasoning, worked autonomously from those prompts. It
searched the OEIS for suitable conjectures, found the proofs, wrote the Lean~4 formalization, and wrote a first draft
of this paper.
\item \emph{Claude Opus~5.5}, used with Claude Code, checked the work independently and then rewrote this paper. It
rebuilt the Lean project from source and confirmed that the theorems listed in Section~\ref{sec:checks} rely only on
Lean's standard axioms; compared the Lean definitions with the OEIS definition; recomputed the sequence with separately
written PARI/GP code, compared it with the terms published in the OEIS, and tested every claim of the paper on many
terms; and read every step of the proofs. It then rewrote the text to make it self-contained and readable. The
theorems and the proof strategies are ChatGPT's. In rewriting, Claude Opus~5.5 expanded the explanations and, in
places, wrote out an argument in more detail or by a more elementary route. The Discussion section was written during
this rewrite; the further observations it contains were checked by computer, and they are not part of the Lean
formalization. Because the result is known, the contribution of this paper is the self-contained
proof and its formal verification.
\end{itemize}
Dev~D.~Goyal has verified the results and is responsible for this paper.

\para{Review status.} This paper has not been peer reviewed, and it has not been submitted to a journal or to the OEIS.

\para{Priority.} An OEIS entry that labels a statement as a conjecture shows only that its contributor did not know
a proof; it does not show that no proof existed. During the research only limited searches for earlier work were
made. No claim of priority is made.

\para{Corrections.} Please report any error in the issue tracker of the repository~\cite{repo}, where the Lean
sources, the numerical checks, and the earlier drafts are also available.

\section*{Acknowledgements}
The sequence comes from the INRIA Encyclopedia of Combinatorial Structures; its present OEIS name is due to Vaclav~Kotesovec and the Stirling-sum formula to Vladimir~Kruchinin. The periodicity conjecture is due to Peter~Bala. The theorem was proved earlier by
Payam~Danesh~\cite{danesh}. In the Lean formalization, the definition of the Stirling numbers and the proof of
Lemma~\ref{lem:power} were adapted from a Mathlib module by Beibei~Xiong, Yu~Shao, Weijie~Jiang and Zhengfeng~Yang, under
the Apache~2.0 license~\cite{stirlingmathlib}.

\begin{thebibliography}{99}\small
\raggedright
\bibitem{oeis} INRIA Encyclopedia of Combinatorial Structures (name by V.~Kotesovec), entry A052895 (Expansion of e.g.f.\ $(1-\sqrt{5-4e^x})/(2(e^x-1))$), \emph{The On-Line Encyclopedia of Integer Sequences}, 2000. \url{https://oeis.org/A052895}. Text as of October 2026.
\bibitem{danesh} P.~Danesh, Eventual periodicity of weighted Stirling transforms and two conjectures of Bala, preprint, version~1, 26 August 2026. \url{https://doi.org/10.33774/coe-2026-4d6jr}.
\bibitem{schreyer} B.~Schreyer, Rigged horse numbers and their modular periodicity, arXiv:2409.03799 (2024).
\bibitem{concrete} R.~L.~Graham, D.~E.~Knuth and O.~Patashnik, \emph{Concrete Mathematics}, second edition, Addison-Wesley, 1994.
\bibitem{stirlingmathlib} B.~Xiong, Y.~Shao, W.~Jiang and Z.~Yang, \texttt{Mathlib/Combinatorics/Enumerative/Stirling.lean}, mathlib4 commit \texttt{4ed598f44360} (2025), Apache License~2.0.
\bibitem{oeisdb} N.~J.~A.~Sloane and The OEIS Foundation Inc., \emph{The On-Line Encyclopedia of Integer Sequences}. \url{https://oeis.org}.
\bibitem{stanley} R.~P.~Stanley, \emph{Enumerative Combinatorics}, Volume~2, Cambridge University Press, 1999. (Chapter~5 treats formal power series, composition and compositional inverses.)
\bibitem{niven} I.~Niven, Formal power series, \emph{American Mathematical Monthly} 76 (1969), 871--889.
\bibitem{lean} L.~de~Moura and S.~Ullrich, The Lean~4 theorem prover and programming language, in \emph{Automated Deduction -- CADE 28}, Lecture Notes in Computer Science 12699, Springer, 2021, 625--635.
\bibitem{mathlib} The mathlib Community, The Lean mathematical library, in \emph{Proceedings of the 9th ACM SIGPLAN International Conference on Certified Programs and Proofs (CPP 2020)}, 367--381. Version used: commit \texttt{c44e0c8ee63c}.
\bibitem{repo} D.~D.~Goyal, \emph{Proofs and corrections for 34 OEIS entries, with Lean~4 formalizations}, repository, 2026. \url{https://github.com/dvgyl/oeis-a396843-proof}.
\end{thebibliography}
\end{document}
